\section{Selecting projections with a prescribed trace}\label{n:sec:selection}

Our first goal is a projection in $D$ with a prescribed trace function
and exact support under a positive contraction. Comparison tests support
ranks, whereas this goal prescribes ordinary traces. The next lemma
connects the two: it approximates a trace by a finite-matrix positive
contraction whose excess support rank is small. After constructing the
projection, we refine the selection by a one-sided cost bound that will
control overlaps and model errors.

\begin{lemma}\label{n:lem:rank}
For $g\in A_+$ and $\xi>0$ there is a positive contraction $k$ in a finite
matrix algebra over $A$ such that, for every $\sigma\in T$,
\begin{equation}\label{n:eq:rank}
 |\sigma(k)-\sigma(g)|<\xi,\qquad
 0\le d_\sigma(k)-\sigma(k)<\xi.
\end{equation}
\end{lemma}
\begin{proof}
First choose a nonzero $b\in A_+$ with $d_\sigma(b)<\xi$ for all
$\sigma\in T$. An extreme point of the compact convex set $T$ has a
factorial tracial GNS representation. This factor cannot be finite
dimensional: the representation is faithful, and would make $A$
finite dimensional. It is therefore a $\mathrm{II}_1$ factor and contains
unital matrices of every size. By \eqref{n:eq:kaplansky} and order-zero
lifting, there is a nonzero order-zero map $M_l\to A$ for every $l$.
The support projections of its diagonal images are orthogonal and
equivalent in every tracial representation. Their sum has trace at most
one, so each diagonal has dimension at most $1/l$. Taking $l>1/\xi$
gives $b$.

Simplicity now gives $c\in\Her(b)_+$ satisfying
$|\sigma(c)-\sigma(g)|<\xi$ on $T$, without a norm bound on $c$.
Here are details of this trace transfer. Approximate $g^{1/2}$ in the
ideal generated by $b$ by $UD_bV$, with $U$ a row, $V$ a column, and
$D_b=\operatorname{diag}(b,\ldots,b)$, so that
$(UD_bV)^*(UD_bV)$ is within $\xi$ of $g$. Set
\[
 C=(D_bU^*UD_b)^{1/2}\in M_N(\Her(b)),\qquad v=CV.
\]
Each entry $v_j$ belongs to $\overline{bA}$. Thus
$c=\sum_jv_jv_j^*\in\Her(b)_+$, and traciality gives
$\sigma(c)=\sigma(v^*v)$, proving the assertion.

For an integer $N\ge\|c\|$, put
\[
 f_j(s)=\min(1,\max(0,s-j+1)),\qquad
 k=\bigoplus_{j=1}^N f_j(c).
\]
On the spectrum of $c$, the functions sum to $s$, while
\[
 0\le\sum_{j=1}^N1_{\{f_j(s)>0\}}-s\le1_{\{s>0\}}.
\]
Hence $\sigma(k)=\sigma(c)$ and
$d_\sigma(k)-\sigma(k)\le d_\sigma(c)\le d_\sigma(b)<\xi$.
\end{proof}

\begin{proposition}\label{n:prop:selection}
Let $h,z,a\in D$ be positive contractions satisfying $hz=z$ and
$\tau(a)\le\tau(z)$ for all $\tau\in \Lambda$. There is a projection $p\in D$
such that
\[
 hp=p,\qquad \tau(p)=\tau(a)\quad(\tau\in \Lambda).
\]
\end{proposition}
\begin{proof}
We first compare finite-matrix approximations to the desired trace with
a spectral cut of the target. A small auxiliary matrix corner supplies
the strict rank inequality needed for comparison. We then discard that
corner and obtain elements of $A$ whose trace and projection defects
vanish along $\omega$.

Take positive contraction lifts $h_n,z_n,a_n$. Set
\[
 \alpha_n=\sup_{\sigma\in T}(\sigma(a_n)-\sigma(z_n))_+,
 \qquad \beta_n=\|(1-h_n)z_n\|_{2,T}.
\]
Both tend to zero along $\omega$. Choose $0<\delta_n<1$ tending to zero
slowly enough that $\beta_n/\delta_n\to_\omega0$, and put
\[
 \gamma_n=\alpha_n+\beta_n/\delta_n+\delta_n,
 \qquad H_n=(h_n-1+\delta_n)_+.
\]
In a tracial representation, the support $Q$ of $H_n$ satisfies
$\delta_n(1-Q)\le1-h_n$. Since $0\le z_n\le1$, tracial
Cauchy--Schwarz gives
\[
 \sigma(z_n)\le\sigma(Q)+\sigma((1-Q)z_n)
 \le d_\sigma(H_n)+\beta_n/\delta_n.
\]
Consequently $\sigma(a_n)\le d_\sigma(H_n)+\gamma_n$.

Apply Lemma~\ref{n:lem:rank} with $g=a_n$ and $\xi=\gamma_n$ to obtain
$k_n$. Also choose a positive contraction $u\in A$ with
$\inf_T\sigma(u)>7/8$, and apply that lemma to $8\gamma_nu$ with the
same tolerance. Its output $c_n$ satisfies
\[
 5\gamma_n\le d_\sigma(c_n)\le10\gamma_n.
\]
Thus, in disjoint finite matrix blocks,
\begin{equation}\label{n:eq:buffer}
 d_\sigma(k_n)\le d_\sigma(H_n)+3\gamma_n
       <d_\sigma(H_n\oplus c_n)\qquad(\sigma\in T).
\end{equation}
The target is nonzero, and all these ranks are finite at each coordinate.
Every densely finite trace is a multiple of a state in $T$; hence the
comparison hypothesis gives $k_n\precsim H_n\oplus c_n$.

We use the cutdown in a form that also controls its trace. Suppress $n$,
and write $b=H_n\oplus c_n$. Choose $x$ with $\|x^*bx-k\|<1/4$ and set
$d=(k-1/2)_+^{1/2}$, $y=b^{1/2}xd$. Then $|y|^2\ge d^2/4$.
For the polar decomposition $y=v|y|$ in the bidual, the support of $d$
lies under $v^*v$, and
\[
 \|(|y|+\epsilon)^{-1}d\|\le2,
 \qquad
 \|vd-y(|y|+\epsilon)^{-1}d\|\le2\epsilon.
\]
The first inequality follows by conjugating $d^2\le4|y|^2$ with
$(|y|+\epsilon)^{-1}$ and taking norms. The second proves that $vd$
belongs to the closed right ideal supported by $b$. Therefore
\[
 w=\sqrt2vd,\qquad w^*w=2(k-1/2)_+,
 \qquad r=ww^*\in\Her(b),\qquad 0\le r\le1.
\]
For $f(s)=2(s-1/2)_+$ on $[0,1]$,
\[
 0\le s-f(s)\le1_{\{s>0\}}-s,
 \qquad 0\le f(s)-f(s)^2\le1_{\{s>0\}}-f(s).
\]
Traciality transfers these estimates from $f(k_n)=w^*w$ to $r_n=ww^*$.
By \eqref{n:eq:rank}, both $|\sigma(r_n)-\sigma(a_n)|$ and
$\sigma(r_n-r_n^2)$ tend uniformly to zero along $\omega$.

Compress $r_n$ to its $H_n$ matrix slot, and identify the resulting
positive contraction $p_n$ with an element of $A$. The complementary
corner is supported by $c_n$, so this removes at most $10\gamma_n$ of
trace. If $P$ is the first-slot projection, the defect identity is
\[
 Pr_nP-(Pr_nP)^2
 =P(r_n-r_n^2)P+Pr_n(1-P)r_nP.
\]
The last term has trace at most $d_\sigma(c_n)$, since $r_n^2\le r_n$.
In particular, both $|\sigma(p_n)-\sigma(a_n)|$ and
$\sigma(p_n-p_n^2)$ are at most $12\gamma_n$, uniformly on $T$.
Moreover $p_n\in\Her(H_n)$, so $\|(1-h_n)p_n\|\le\delta_n$.
The class $p=[(p_n)]$ is a projection: its positive contraction defects
have vanishing trace, hence vanishing uniform $2$-norm. It has all the
required properties.
\end{proof}

\begin{lemma}\label{n:lem:cost}
Let $q\in D$ be a projection, $Y\in(qDq)_+$, $0<t<1$, and
$\varepsilon>0$. There is a projection $r\le q$ such that, for $\tau\in \Lambda$,
\begin{equation}\label{n:eq:cost-corner}
 \tau(r)=t\tau(q),\qquad
 \tau(rY)\le t\bigl(\tau(Y)+\varepsilon\tau(q)\bigr).
\end{equation}
\end{lemma}
\begin{proof}
If $q=0$, take $r=0$; if $Y=0$, apply Proposition~\ref{n:prop:selection}
to $(h,z,a)=(q,q,tq)$. Scaling $Y$ and the tolerance reduces to
$0\le Y\le q$. All functions below are evaluated at $Y$ in the unital
corner $qDq$.

Choose a large integer $m$. Let $F_0=0$, $F_m=1$; for $1\le j<m$ let
$F_j$ be $1$ up to $(j-1/2)/m$, decrease linearly to $0$ at $j/m$, and
vanish thereafter. Let $G_j$ be $1$ up to $j/m$, decrease to $0$ by
$(j+1/3)/m$, and vanish thereafter for $j<m$; put $G_m=1$. These functions
satisfy
\[
 G_jF_j=F_j,\qquad F_{j+1}G_j=G_j.
\]
Construct orthogonal projections $r_j$ recursively with
\[
 G_jr_j=r_j,\qquad \tau(r_j)=t\tau(F_j-F_{j-1}).
\]
Indeed, if $R=\sum_{i<j}r_i$, then $F_jR=G_jR=R$ and
$\tau(R)=t\tau(F_{j-1})$. Apply Proposition~\ref{n:prop:selection} to
\[
 h=G_j-R,\qquad z=F_j-R,\qquad a=t(F_j-F_{j-1}).
\]
They are positive contractions, $hz=z$, and
$\tau(z)-\tau(a)=(1-t)\tau(F_j)\ge0$. The selected projection is
orthogonal to $R$, lies under $q$, and is fixed by $G_j$.

Now $r=\sum_jr_j$ has trace $t\tau(q)$. The spectral support conditions
give
\[
 r_jYr_j\le\frac{j+1/3}{m}r_j,
\]
whereas $F_j-F_{j-1}$ is supported at values at least $(j-3/2)/m$.
The difference of these thresholds is $11/(6m)<2/m$. Summing therefore
gives
\[
 \tau(rY)\le t\sum_j\frac{j+1/3}{m}\tau(F_j-F_{j-1})
       \le t\bigl(\tau(Y)+(2/m)\tau(q)\bigr).
\]
Take $2/m<\varepsilon$.
\end{proof}

\begin{corollary}\label{n:cor:cost}
If $E\subset D$ is finite dimensional and unital, $Y\in D_+$,
$0<t<1$, and $\varepsilon>0$, there is a projection $p\in D\cap E'$ with
\begin{equation}\label{n:eq:cost}
 \tau(pd)=t\tau(d)\quad(d\in E),\qquad
 \tau(pY)\le t\bigl(\tau(Y)+\varepsilon\bigr)\quad(\tau\in \Lambda).
\end{equation}
\end{corollary}
\begin{proof}
For a matrix summand of size $l$ with matrix units $e_{ij}$, apply
Lemma~\ref{n:lem:cost} in $e_{11}De_{11}$ to the observable
$l^{-1}\sum_j e_{1j}Ye_{j1}$. Repeat the resulting projection along the
diagonal, using the matrix units. The repeated projection commutes with
the summand and splits its traces proportionally. Its cost is $l$ times
the selected corner cost. Adding over the orthogonal summands gives
\eqref{n:eq:cost}, since their units sum to $1_D$.
\end{proof}
